% !TEX program = xelatex
\documentclass[11pt,a4paper]{article}

% ── SCX Framework Standard Packages ────────────────────────────────────
\usepackage{fontspec}
\usepackage{amsmath,amssymb,amsthm}
\usepackage{mathtools}
\usepackage{bm}
\usepackage{graphicx}
\usepackage{xcolor}
\usepackage{hyperref}
\usepackage{geometry}
\usepackage{enumitem}
\usepackage{booktabs}
\usepackage{multirow}
\usepackage{array}
\usepackage{algorithm}
\usepackage{algpseudocode}
\usepackage{setspace}
\usepackage{dsfont}
\usepackage{bbm}
\usepackage{tikz}

% CJK support via fontspec (XeLaTeX/LuaLaTeX)
\usepackage{xeCJK}
\setCJKmainfont{SimSun}

\geometry{margin=1in}
\onehalfspacing

% ── SCX Theorems & Proofs ──────────────────────────────────────────────
\newtheorem{theorem}{Theorem}[section]
\newtheorem{proposition}[theorem]{Proposition}
\newtheorem{lemma}[theorem]{Lemma}
\newtheorem{corollary}[theorem]{Corollary}
\newtheorem{conjecture}[theorem]{Conjecture}
\newtheorem{definition}[theorem]{Definition}
\newtheorem{remark}[theorem]{Remark}
\theoremstyle{definition}
\newtheorem{assumption}[theorem]{Assumption}
\newtheorem{example}[theorem]{Example}

% ── Rigor Labels ───────────────────────────────────────────────────────
\definecolor{rigorcolor}{RGB}{0,120,0}
\newcommand{\rigorFull}{{\color{rigorcolor}$\blacksquare$\,\textbf{[Full Proof]}}}

% ── SCX Framework Macros ───────────────────────────────────────────────
\newcommand{\Situs}{\textsc{Situs}}
\newcommand{\Yajie}{\textsc{Yajie}}
\newcommand{\Spring}{\textsc{Spring}}
\newcommand{\Cercis}{\textsc{Cercis}}
\newcommand{\SCX}{\textsc{SCX}}

% ── Math Operators ─────────────────────────────────────────────────────
\DeclareMathOperator*{\argmin}{arg\,min}
\DeclareMathOperator*{\argmax}{arg\,max}
\DeclareMathOperator{\TV}{TV}
\DeclareMathOperator{\KL}{KL}
\DeclareMathOperator{\Var}{Var}
\DeclareMathOperator{\Cov}{Cov}
\DeclareMathOperator{\Bias}{Bias}
\DeclareMathOperator{\Corr}{Corr}
\DeclareMathOperator{\diag}{diag}
\DeclareMathOperator{\supp}{supp}

% ── Notation ───────────────────────────────────────────────────────────
\newcommand{\R}{\mathbb{R}}
\newcommand{\N}{\mathbb{N}}
\newcommand{\E}{\mathbb{E}}
\newcommand{\Pbb}{\mathbb{P}}
\newcommand{\ind}[1]{\mathbf{1}[#1]}
\newcommand{\norm}[1]{\left\lVert#1\right\rVert}
\newcommand{\abs}[1]{\left|#1\right|}
\newcommand{\inner}[2]{\langle #1, #2 \rangle}
\newcommand{\cF}{\mathcal{F}}
\newcommand{\cL}{\mathcal{L}}
\newcommand{\cM}{\mathcal{M}}
\newcommand{\cP}{\mathcal{P}}
\newcommand{\cS}{\mathcal{S}}
\newcommand{\cT}{\mathcal{T}}
\newcommand{\cX}{\mathcal{X}}

% ── NV-Center-specific macros ──────────────────────────────────────────
\newcommand{\Ttwo}{T_2}
\newcommand{\TtwoStar}{T_2^*}
\newcommand{\Tone}{T_1}
\newcommand{\NV}{NV}
\newcommand{\NVminus}{NV${}^-$}
\newcommand{\ODMR}{ODMR}
\newcommand{\CVD}{CVD}
\newcommand{\HPHT}{HPHT}
\newcommand{\RabiFreq}{\Omega_R}
\newcommand{\FidReadout}{\mathcal{F}_{\text{ro}}}
\newcommand{\DiamondPurity}{\mathcal{P}_{\text{dia}}}
\newcommand{\SurfNoise}{\mathcal{N}_{\text{surf}}}
\newcommand{\BulkImp}{\mathcal{I}_{\text{bulk}}}
\newcommand{\ProtocolNoise}{\mathcal{N}_{\text{prot}}}

% ── Assumption / Limitation Tags ───────────────────────────────────────
\newcommand{\assumptionTag}[1]{\textsf{\textbf{[A#1]}}}
\newcommand{\limitationTag}[1]{\textsf{\textbf{[L#1]}}}

% ════════════════════════════════════════════════════════════════════════
% TITLE
% ════════════════════════════════════════════════════════════════════════
\title{\textbf{SCX-Audited NV Centers:}\\[4pt]
       \large Multi-Laboratory Certification of Diamond Qubit Quality Claims\\[4pt]
       \small SCX审计氮空位中心：金刚石量子比特质量声明的多实验室认证}
\author{SCX Research Collective \\ \small Xiaogan Supercomputing Center}
\date{\today}

\begin{document}
\maketitle

% ════════════════════════════════════════════════════════════════════════
% ABSTRACT
% ════════════════════════════════════════════════════════════════════════
\begin{abstract}
\noindent
{\bf 摘要.} 氮空位（\NV{}）中心是金刚石中的原子尺度缺陷，构成固态量子传感、量子网络与量子信息处理的核心平台。然而，\NV{}中心量子比特的品质声明——相干时间（相干时间）、拉比频率（拉比频率）、读取保真度（读取保真度）——在不同实验室之间存在显著差异，其根源包括金刚石纯度、表面处理、应变分布及测量协议的多样性。本文将多实验室\NV{}中心审计形式化为\SCX{}多专家认证问题：每个\NV{}中心样品为一个状态（state），每个实验室为一名专家（expert），不同的测量协议（\ODMR{}、Ramsey干涉、Hahn回波、动力学去耦）构成多模态专家（multi-modal experts），而金刚石纯度/表面/应变变异构成状态异质性（state heterogeneity）。

本文发展三个核心定理并给出\rigorFull{}完整证明：
(1)~{\bf 多实验室误差检测定理}：$M$个独立实验室审计同一\NV{}属性（如$\Ttwo$）时，所有实验室同时遗漏系统性测量误差的概率受$\exp(-2M_{\text{eff}}\,\Delta^2)$约束，其中$M_{\text{eff}} = M/(1+(M-1)\bar{\rho})$修正实验室间因共享校准标准、相近设备或共同协议基准产生的相关性；
(2)~{\bf 跨协议Cercis评分收敛定理}：在弱依赖条件下，\Cercis{}评分$S = Q + \eta N$（其中$Q$聚合相干时间再现性与读取保真度一致性，$N$量化金刚石制备方法新颖性）收敛至真实品质度量，其收敛速率受$M_{\text{eff}}$与测量协议多样性的联合约束；
(3)~{\bf $\Ttwo$退化源不可辨识定理}：当实验室报告不同$\Ttwo$值时，退化来源（表面噪声 vs 体内杂质 vs 测量协议差异）在未声明结构假设的前提下不可辨识——通过构造观测等价的三个世界给出证明。

进一步引入{\bf \Cercis{}评分} $S = Q + \eta N$，其中$Q$聚合相干时间再现性与读取保真度一致性，$N$量化金刚石制备方法新颖性（基于CVD/HPHT衬底、N-14/N-15同位素、浅层/深层\NV{}植入深度）。{\bf \Yajie{}多实验室共识}以\Cercis{}评分为权重聚合$M$个实验室的测量结果。实验基准覆盖：N-14 vs N-15同位素对比、浅层（$<10$~nm）vs 深层（$>50$~nm）\NV{}中心、CVD vs HPHT金刚石衬底、以及四种测量协议（\ODMR{}、Ramsey、Hahn回波、XY8动力学去耦）的交叉比较。

本文不做综述，不宣称量子霸权。本文提供定理与证明。
\end{abstract}

\vspace{4pt}
\noindent{\small\bfseries Keywords:}
NV center, diamond qubit, coherence time, quantum sensing, multi-laboratory audit,
error detection, Hoeffding bound, unidentifiability, Cercis score, Yajie consensus,
ODMR, Ramsey interferometry, Hahn echo, dynamical decoupling, CVD, HPHT.

\vspace{4pt}
\noindent{\small\bfseries 关键词:}
氮空位中心, 金刚石, 相干时间, 量子传感, 多实验室审计, 误差检测,
Hoeffding界, 不可辨识性, Cercis评分, Yajie共识, ODMR, Ramsey干涉, Hahn回波, 动力学去耦.
\vspace{8pt}

% ════════════════════════════════════════════════════════════════════════
\section{Introduction 引言}
\label{sec:intro}
% ════════════════════════════════════════════════════════════════════════

The nitrogen-vacancy (\NV{}) center in diamond is among the most intensively studied solid-state quantum systems. A single \NVminus{} center—a substitutional nitrogen atom adjacent to a lattice vacancy—possesses a spin-triplet ground state with optical initialization and readout, millisecond-scale coherence times at room temperature~\cite{doherty2013nitrogen,barry2020sensitivity}, and coherent coupling to proximal nuclear spins~\cite{bradley2019ten}. These properties make \NV{} centers a leading platform for quantum sensing~\cite{staudacher2013nuclear}, quantum networks~\cite{hermans2023entangling}, and quantum information processing~\cite{abrahao2025quantum}.

Yet the field grapples with a persistent and under-theorized problem: \textbf{quality claims about \NV{} centers are not systematically auditable across laboratories.} A research group at Laboratory A reports $\Ttwo = 1.2$~ms for a shallow \NV{} center in CVD diamond; Laboratory B reports $\Ttwo = 180$~$\mu$s for a nominally identical sample; Laboratory C measures $\Ttwo = 2.4$~ms using a dynamical decoupling sequence that Laboratory A did not employ. The reader, the reviewer, and the funding agency have no principled framework for determining whether these discrepancies reflect genuine sample heterogeneity, measurement protocol differences, calibration errors, or some combination thereof.

\medskip
\noindent\textbf{The Audit Gap in Quantum Materials.} The current publication model for \NV{} center characterization operates as follows:
\begin{enumerate}[label=(\roman*)]
    \item A single laboratory fabricates or procures a diamond sample;
    \item The laboratory performs measurements using its chosen protocols and equipment;
    \item The laboratory reports the measured parameters ($\Ttwo$, $\Tone$, $\RabiFreq$, $\FidReadout$) as intrinsic properties of the \NV{} center;
    \item Peer review checks for internal consistency, methodological correctness, and novelty of claims — but \textbf{does not require independent replication of the measurements by a different laboratory.}
\end{enumerate}
This creates an \textit{audit gap}: the measured coherence time of an \NV{} center is simultaneously a property of the diamond, a property of the measurement protocol, and a property of the laboratory's calibration. Without a structured framework for multi-laboratory certification, the community cannot distinguish between a genuinely superior diamond sample and a measurement artifact.

\medskip
\noindent\textbf{The SCX Solution.} The \SCX{} (Supercomputing Certification eXchange) framework~\cite{scx2025} treats quality certification as a multi-expert consensus problem. In the context of \NV{} centers:
\begin{itemize}[nosep]
    \item Each \textbf{\NV{} center sample} is a \textbf{state} $s \in \cS$, characterized by diamond substrate (CVD or HPHT), nitrogen isotope (N-14 or N-15), implantation depth, surface termination, and local strain environment;
    \item Each \textbf{laboratory} is an \textbf{expert} $\ell \in \cL$, equipped with measurement apparatus, calibration standards, and protocol preferences;
    \item Each \textbf{measurement protocol} ($\ODMR{}$, Ramsey interferometry, Hahn echo, XY$n$ dynamical decoupling) is a \textbf{modality}, providing a distinct projection of the underlying qubit quality;
    \item \textbf{Heterogeneity} in diamond purity, surface noise, and strain distribution constitutes the state space over which quality claims must be certified.
\end{itemize}

\medskip
\noindent\textbf{This Paper.} We formalize multi-laboratory \NV{} center auditing as an \SCX{} multi-expert certification problem. We prove three theorems with full mathematical rigor ($\rigorFull$):
\begin{enumerate}[label=(\roman*)]
    \item \textbf{Theorem~1 (Multi-Laboratory Error Detection)}: The probability that $M$ laboratories all miss a systematic measurement error in an \NV{} property decays exponentially in the effective number of independent laboratories $M_{\text{eff}}$, with bound $\exp(-2M_{\text{eff}}\,\Delta^2)$;
    \item \textbf{Theorem~2 (Cross-Protocol \Cercis{} Score Convergence)}: The \Cercis{} score $S = Q + \eta N$, which aggregates coherence time reproducibility ($Q$) and fabrication method novelty ($N$), converges to a true quality metric under weak dependence conditions;
    \item \textbf{Theorem~3 ($\Ttwo$ Degradation Source Unidentifiability)}: When laboratories report discrepant $\Ttwo$ values, the causal source—surface noise, bulk paramagnetic impurities, or measurement protocol differences—is fundamentally unidentifiable without declared structural assumptions.
\end{enumerate}

We specify:
\begin{itemize}[nosep]
    \item The \textbf{\Cercis{} Score} $S = Q + \eta N$, where $Q$ aggregates coherence time reproducibility across laboratories and readout fidelity consistency across protocols, and $N$ quantifies the novelty of the diamond fabrication method;
    \item The \textbf{\Yajie{} Multi-Laboratory Consensus} procedure, which weights each laboratory's measurement by its \Cercis{} score on calibration samples;
    \item An \textbf{experimental benchmark protocol} covering four comparison axes: N-14 vs N-15 isotopic composition, shallow ($<10$~nm) vs deep ($>50$~nm) \NV{} implantation, CVD vs HPHT diamond substrates, and four measurement protocols (\ODMR{}, Ramsey, Hahn echo, XY8 dynamical decoupling).
\end{itemize}

\textbf{What this paper is not.} This is not a review of \NV{} center physics. It is not an experimental demonstration. It does not claim quantum supremacy. It is a mathematical paper about multi-laboratory consensus, error detection probability bounds, unidentifiability, and auditable quality metrics for diamond quantum devices.

% ════════════════════════════════════════════════════════════════════════
\section{Formalization 形式化}
\label{sec:formalism}
% ════════════════════════════════════════════════════════════════════════

\subsection{NV Center State Space}

\begin{definition}[NV Center State 氮空位中心状态]
\label{def:nv_state}
An \NV{} center sample is characterized by a state vector
$\mathbf{s} = (s_1, \dots, s_d) \in \cS \subset \R^d$, where $d$ is the
dimensionality of the relevant physical parameters. The state components include:
\begin{itemize}[nosep]
    \item \textbf{Diamond substrate type}: $s_{\text{sub}} \in \{\CVD{}, \HPHT{}\}$,
    indicating chemical vapor deposition or high-pressure high-temperature synthesis;
    \item \textbf{Nitrogen isotope}: $s_{\text{iso}} \in \{\text{N-14}, \text{N-15}\}$,
    where N-15 provides nuclear spin $I=1/2$ (vs $I=1$ for N-14) and
    different hyperfine coupling to the \NV{} electronic spin;
    \item \textbf{Implantation depth}: $s_{\text{depth}} \in \R_+$, the distance
    from the diamond surface to the \NV{} center (shallow: $<10$~nm;
    deep: $>50$~nm; intermediate: $10$--$50$~nm);
    \item \textbf{Surface termination}: $s_{\text{surf}} \in \{\text{H}, \text{O}, \text{F}, \text{N}\}$,
    chemical species terminating the diamond surface, affecting surface noise;
    \item \textbf{Nitrogen concentration}: $s_{\text{[N]}} \in \R_+$, the
    bulk nitrogen impurity concentration (ppm), controlling paramagnetic spin bath density;
    \item \textbf{Local strain}: $s_{\text{strain}} \in \R^6$, the local
    strain tensor at the \NV{} site, affecting zero-field splitting and
    optical properties;
    \item \textbf{${}^{13}$C isotopic abundance}: $s_{\text{C13}} \in [0, 1]$,
    natural (1.1\%) or isotopically purified ($<0.01\%$), controlling nuclear spin bath density.
\end{itemize}
The state space is $\cS = \{\CVD, \HPHT\} \times \{\text{N-14}, \text{N-15}\}
\times \R_+ \times \{\text{H}, \text{O}, \text{F}, \text{N}\} \times \R_+ \times \R^6 \times [0,1]$.
\end{definition}

\begin{definition}[True Quality Parameters 真实品质参数]
\label{def:quality}
For an \NV{} center in state $\mathbf{s}$, the true (but unobservable) quality
parameters are:
\begin{align}
    \Ttwo(\mathbf{s}) &\in \R_+ \quad \text{(inhomogeneous dephasing time)}, \\
    \Tone(\mathbf{s}) &\in \R_+ \quad \text{(longitudinal relaxation time)}, \\
    \RabiFreq(\mathbf{s}) &\in \R_+ \quad \text{(Rabi frequency, controllable via MW power)}, \\
    \FidReadout(\mathbf{s}) &\in [0,1] \quad \text{(single-shot readout fidelity)}.
\end{align}
These parameters are deterministic functions of the state $\mathbf{s}$,
i.e., they are intrinsic physical properties. The measurement challenge
is that no laboratory can access $\Ttwo(\mathbf{s})$ directly; each
laboratory observes an estimate $\widehat{T}_2^{(\ell, p)}(\mathbf{s})$
that depends on both the laboratory $\ell$ and the measurement protocol $p$.
\end{definition}

\subsection{Laboratories as Experts}

\begin{definition}[Laboratory Expert 实验室专家]
\label{def:lab}
A laboratory $\ell \in \cL = \{1, \dots, L\}$ is characterized by:
\begin{itemize}[nosep]
    \item \textbf{Measurement apparatus}: optical setup (confocal microscope,
    NA, laser wavelength, detector type), microwave delivery (antenna/waveguide
    design, amplifier specifications), magnetic field control (Helmholtz coils,
    permanent magnets, shielding);
    \item \textbf{Calibration standards}: reference samples, calibration
    protocols, data analysis pipelines;
    \item \textbf{Protocol repertoire}: $\cP_\ell \subseteq \cP$, the set of
    measurement protocols available at laboratory $\ell$;
    \item \textbf{Systematic bias}: $\beta_\ell(p, \mathbf{s}) \in \R$, the
    systematic measurement bias of laboratory $\ell$ using protocol $p$ on
    state $\mathbf{s}$.
\end{itemize}
\end{definition}

\begin{definition}[Measurement Protocols 测量协议]
\label{def:protocols}
The set of measurement protocols $\cP$ comprises:
\begin{enumerate}[label=\textbf{P\arabic*:}, nosep]
    \item \textbf{ODMR (optically detected magnetic resonance)}:
    Continuous-wave or pulsed measurement of the ground-state spin resonance,
    yielding $\TtwoStar(\mathbf{s})$ (inhomogeneous dephasing including
    quasi-static bath fluctuations);
    \item \textbf{Ramsey interferometry}: Free-induction decay measurement
    using $\pi/2$--$\tau$--$\pi/2$ pulse sequence, sensitive to low-frequency
    noise, yielding $\TtwoStar(\mathbf{s})$ with different spectral filter
    function;
    \item \textbf{Hahn echo}: $\pi/2$--$\tau$--$\pi$--$\tau$--$\pi/2$ sequence,
    refocusing static inhomogeneity, yielding $\Ttwo(\mathbf{s})$;
    \item \textbf{Dynamical decoupling (XY$n$)}: Periodic $\pi$-pulse trains
    ($\text{XY}4$, $\text{XY}8$, $\text{XY}16$), extending coherence to
    $T_{2,\text{DD}} > \Ttwo(\mathbf{s})$ by filtering specific noise
    frequency bands.
\end{enumerate}
Each protocol $p \in \cP$ has an associated filter function
$F_p(\omega): \R_+ \to [0,1]$ describing its spectral sensitivity to
magnetic noise at frequency $\omega$. Different protocols are sensitive
to different noise spectral regions, making them \textbf{multi-modal experts}
that provide complementary information about the decoherence environment.
\end{definition}

\begin{definition}[Measurement Model 测量模型]
\label{def:measurement}
For laboratory $\ell$ measuring quality parameter $Q \in \{\Ttwo, \Tone, \RabiFreq, \FidReadout\}$
on state $\mathbf{s}$ using protocol $p$, the observed value is:
\begin{equation}
    \widehat{Q}^{(\ell, p)}(\mathbf{s}) = Q(\mathbf{s}) + \beta_\ell(p, \mathbf{s}) + \varepsilon_{\ell, p, \mathbf{s}},
    \label{eq:measurement_model}
\end{equation}
where:
\begin{itemize}[nosep]
    \item $Q(\mathbf{s})$ is the true (unknown) parameter value;
    \item $\beta_\ell(p, \mathbf{s})$ is the systematic bias of laboratory $\ell$
    using protocol $p$ on state $\mathbf{s}$;
    \item $\varepsilon_{\ell, p, \mathbf{s}} \sim \cN(0, \sigma_{\ell, p}^2(\mathbf{s}))$
    is zero-mean Gaussian measurement noise with state-dependent variance.
\end{itemize}
\end{definition}

\subsection{State Heterogeneity}

\begin{definition}[State Heterogeneity Measure 状态异质性度量]
\label{def:heterogeneity}
The heterogeneity between two \NV{} center states $\mathbf{s}, \mathbf{s}' \in \cS$ is:
\begin{equation}
    d_{\cS}(\mathbf{s}, \mathbf{s}') = \sum_{k=1}^{d} w_k \cdot \delta_k(s_k, s'_k),
    \label{eq:heterogeneity}
\end{equation}
where $\delta_k$ is a domain-appropriate distance (Hamming distance for categorical
variables, normalized absolute difference for continuous variables) and
$w_k \geq 0$ are importance weights with $\sum_k w_k = 1$. State heterogeneity
quantifies the degree to which two \NV{} centers differ in their fabrication
parameters and physical environment.
\end{definition}

\begin{remark}[Sources of State Heterogeneity]
State heterogeneity in \NV{} centers arises from multiple physical mechanisms:
\begin{enumerate}[nosep]
    \item \textbf{Diamond purity}: CVD diamond typically contains fewer
    paramagnetic nitrogen impurities (P1 centers) than HPHT diamond,
    leading to longer $\Ttwo$ in the bulk limit;
    \item \textbf{Surface noise}: Shallow \NV{} centers ($<10$~nm depth)
    experience surface spin noise from dangling bonds and adsorbates,
    dominant over bulk impurity noise for depths below $\sim 5$--$10$~nm;
    \item \textbf{Strain distribution}: Local strain inhomogeneity splits
    and broadens the \NV{} optical and spin transitions, reducing both
    $\TtwoStar$ and readout fidelity;
    \item \textbf{Isotopic composition}: ${}^{13}$C nuclear spins (natural
    abundance $1.1\%$) form a fluctuating spin bath; isotopic purification
    to $<0.01\%$ ${}^{13}$C can extend $\Ttwo$ by an order of magnitude.
\end{enumerate}
These mechanisms are \textit{not independent}: surface treatment affects strain,
nitrogen concentration correlates with substrate choice, and implantation
depth interacts with surface termination chemistry. This coupling makes
individual causal attribution difficult—a fact formalized in Theorem~3.
\end{remark}

% ════════════════════════════════════════════════════════════════════════
\section{Assumptions 假设}
\label{sec:assumptions}
% ════════════════════════════════════════════════════════════════════════

All theorems hold under the following assumptions. Each is stated,
labeled, and discussed with its physical justification and verifiability.

\begin{assumption}[Bounded Measurement Error — \assumptionTag{A1}]
\label{ass:A1}
For every laboratory $\ell \in \cL$, protocol $p \in \cP$, and state
$\mathbf{s} \in \cS$, the measurement error on any quality parameter
$Q \in \{\Ttwo, \Tone, \RabiFreq, \FidReadout\}$ is bounded:
\begin{equation}
    |\widehat{Q}^{(\ell, p)}(\mathbf{s}) - Q(\mathbf{s})| \leq B_Q,
\end{equation}
where $B_Q > 0$ is a known upper bound. For coherence time measurements,
$B_{\Ttwo}$ is typically limited by the maximum pulse sequence duration
and laser power stability.
\end{assumption}

\begin{assumption}[Laboratory Independence Given State — \assumptionTag{A2}]
\label{ass:A2}
For laboratories $\ell \neq \ell'$ with disjoint equipment,
independent calibration, and no shared samples, their measurement
errors are conditionally independent given the state $\mathbf{s}$:
\begin{equation}
    \varepsilon_{\ell, p, \mathbf{s}} \perp\!\!\!\perp \varepsilon_{\ell', p', \mathbf{s}} \mid \mathbf{s}.
\end{equation}
For laboratories sharing calibration standards, equipment manufacturers,
or data analysis pipelines, correlation is bounded by the overlap
coefficient $\rho_{\ell\ell'} \leq J_{\ell\ell'} \cdot \gamma$,
where $J_{\ell\ell'}$ is a Jaccard-like overlap index of shared
resources and $\gamma \in [0,1]$ is an attenuation factor.
\end{assumption}

\begin{assumption}[Positive Individual Laboratory Margins — \assumptionTag{A3}]
\label{ass:A3}
For any laboratory $\ell$ and its best protocol $p^* \in \cP_\ell$,
the expected absolute scaled error satisfies:
\begin{equation}
    \E\left[\frac{|\widehat{Q}^{(\ell, p^*)}(\mathbf{s}) - Q(\mathbf{s})|}{B_Q}\right]
    \leq \frac{1}{2} - \Delta_\ell,
\end{equation}
with $\Delta_\ell > 0$. That is, each laboratory's measurement is
strictly better than random guessing within the error bound. For
well-established \NV{} characterization labs, $\Delta_\ell \gtrsim 0.15$--$0.3$.
\end{assumption}

\begin{assumption}[Protocol Diversity — \assumptionTag{A4}]
\label{ass:A4}
The set of protocols $\cP = \{\ODMR, \text{Ramsey}, \text{Hahn}, \text{XY}n\}$
has distinct filter functions $F_p(\omega)$, such that no protocol's
noise sensitivity is a linear combination of the others:
\begin{equation}
    \nexists\, \boldsymbol{\alpha} \in \R^{|\cP|-1} \text{ s.t. } F_p(\omega) = \sum_{q \neq p} \alpha_q F_q(\omega) \; \forall \omega.
\end{equation}
This ensures that protocol diversity provides genuine information gain
rather than redundant noise filtering.
\end{assumption}

\begin{assumption}[Calibration Sample Availability — \assumptionTag{A5}]
\label{ass:A5}
A calibration set $\cS_{\text{cal}} \subset \cS$ of $n_{\text{cal}} \geq 20$
\NV{} center samples with peer-reviewed consensus quality parameters
is available. These samples span the state heterogeneity space
$\cS$ and have $\Ttwo$ values confirmed by at least 3 independent
laboratories using at least 2 distinct protocols each.
\end{assumption}

\begin{assumption}[Finite State Heterogeneity — \assumptionTag{A6}]
\label{ass:A6}
The state heterogeneity metric $d_{\cS}(\mathbf{s}, \mathbf{s}')$ is
bounded: $\sup_{\mathbf{s}, \mathbf{s}' \in \cS} d_{\cS}(\mathbf{s}, \mathbf{s}') \leq D_{\max} < \infty$.
The quality parameters $Q(\mathbf{s})$ are $L_Q$-Lipschitz with respect
to $d_{\cS}$:
\begin{equation}
    |Q(\mathbf{s}) - Q(\mathbf{s}')| \leq L_Q \cdot d_{\cS}(\mathbf{s}, \mathbf{s}').
\end{equation}
\end{assumption}

\begin{assumption}[Unbiasedness of Best Protocol — \assumptionTag{A7}]
\label{ass:A7}
For each laboratory $\ell$, there exists at least one protocol
$p^*_\ell \in \cP_\ell$ such that the systematic bias vanishes
asymptotically with repetition:
\begin{equation}
    \lim_{N_{\text{rep}} \to \infty} \E[\beta_\ell(p^*_\ell, \mathbf{s}) \mid N_{\text{rep}} \text{ repetitions}] = 0.
\end{equation}
For finite repetitions, $|\beta_\ell(p^*_\ell, \mathbf{s})| \leq b_\ell / \sqrt{N_{\text{rep}}}$.
\end{assumption}

\begin{assumption}[Protocol-Independent Ground Truth — \assumptionTag{A8}]
\label{ass:A8}
There exists a true, protocol-independent value of each quality
parameter $Q(\mathbf{s})$. While no single protocol measures this
true value directly (each protocol is sensitive to different noise
spectral bands), the parameter is well-defined as the limit of
increasingly refined dynamical decoupling: $Q(\mathbf{s}) = \lim_{n \to \infty} Q(\text{XY}n, \mathbf{s})$
for coherence times, or the quantum limit for readout fidelity.
\end{assumption}

\begin{remark}[Verifiability of Assumptions]
\ref{ass:A1} is verified by the physical limits of the apparatus
(maximum $\pi$-pulse fidelity, laser stability). \ref{ass:A2} is
verified by documentation of shared resources; independence is the
default under disjoint procurement. \ref{ass:A3} is verified by
each laboratory's track record on calibration samples. \ref{ass:A4}
is verified by the known filter functions of standard pulse sequences.
\ref{ass:A5} requires community curation of a shared calibration
sample set. \ref{ass:A6} is verified by the finite parameter ranges
of diamond synthesis. \ref{ass:A7} is verified by convergence of
repeated measurements. \ref{ass:A8} is a definitional assumption
about the existence of asymptotic limits.
\end{remark}

% ════════════════════════════════════════════════════════════════════════
\section{Theorem 1: Multi-Laboratory Error Detection 多实验室误差检测}
\label{sec:thm1}
% ════════════════════════════════════════════════════════════════════════

\begin{theorem}[Multi-Laboratory Systematic Error Detection Bound 多实验室系统误差检测界限]
\label{thm:error_detection}
\rigorFull

Let $\cL = \{\ell_1, \dots, \ell_L\}$ be a set of $L$ laboratories
auditing the same \NV{} quality parameter $Q$ on state $\mathbf{s}$,
each using their best protocol $p^*_\ell$, satisfying
Assumptions~\ref{ass:A1}--\ref{ass:A3}.

Define the effective number of independent laboratories:
\begin{equation}
    L_{\eff} = \frac{L}{1 + (L-1)\bar{\rho}},
    \label{eq:L_eff}
\end{equation}
where $\bar{\rho} = \frac{1}{L(L-1)} \sum_{\ell \neq \ell'} \rho_{\ell\ell'}$
is the average pairwise correlation of measurement errors, with
$\rho_{\ell\ell'} = \Corr(\varepsilon_{\ell, p^*_\ell, \mathbf{s}},
\varepsilon_{\ell', p^*_{\ell'}, \mathbf{s}})$ bounded by
Assumption~\ref{ass:A2}.

Let $\Delta = \min_\ell \Delta_\ell > 0$ be the minimum laboratory
margin from Assumption~\ref{ass:A3}. Consider the multi-laboratory
consensus estimator:
\begin{equation}
    \widehat{Q}_{\text{consensus}}(\mathbf{s}) = \frac{1}{L} \sum_{\ell=1}^{L} \widehat{Q}^{(\ell, p^*_\ell)}(\mathbf{s}).
    \label{eq:consensus}
\end{equation}

Then the probability that the consensus deviates from the true value
by more than $\delta > 0$ after normalization is bounded by:
\begin{equation}
    \Pbb\left(\frac{|\widehat{Q}_{\text{consensus}}(\mathbf{s}) - Q(\mathbf{s})|}{B_Q} > \delta \right)
    \leq 2 \exp\!\bigl(-2 L_{\eff} \, \max(\Delta - \delta, 0)^2 \bigr).
    \label{eq:hoeffding_bound}
\end{equation}

In particular, for the systematic error detection event — that all
$L$ laboratories simultaneously produce measurements consistent with
a spurious value $\widetilde{Q} \neq Q(\mathbf{s})$ such that
$|\widetilde{Q} - Q(\mathbf{s})| > \delta B_Q$ — the probability satisfies:
\begin{equation}
    \Pbb(\text{all } L \text{ labs miss systematic error of magnitude } > \delta B_Q)
    \leq \exp\!\bigl(-2 L_{\eff} \, \Delta^2 \bigr).
    \label{eq:systematic_bound}
\end{equation}
\end{theorem}

\begin{proof}
The proof proceeds in five steps, adapting the Azuma--Hoeffding
inequality to account for inter-laboratory correlation.

\medskip\noindent\textit{Step 1: Normalization and centering.}
For each laboratory $\ell$, define the normalized error:
\begin{equation}
    Z_\ell = 1 - \frac{|\widehat{Q}^{(\ell, p^*_\ell)}(\mathbf{s}) - Q(\mathbf{s})|}{B_Q}.
\end{equation}
By Assumption~\ref{ass:A1}, $|\widehat{Q}^{(\ell, p^*_\ell)} - Q| \leq B_Q$,
so $Z_\ell \in [0,1]$. By Assumption~\ref{ass:A3},
\begin{equation}
    \E[Z_\ell] = 1 - \frac{\E[|\widehat{Q}^{(\ell, p^*_\ell)} - Q|]}{B_Q}
              \geq 1 - \left(\frac{1}{2} - \Delta_\ell\right)
              = \frac{1}{2} + \Delta_\ell
              \geq \frac{1}{2} + \Delta.
\end{equation}
Thus each $Z_\ell$ is a random variable bounded in $[0,1]$ with
expectation at least $1/2 + \Delta$.

\medskip\noindent\textit{Step 2: Effective sample size from inter-laboratory correlation.}
The $L$ laboratories have measurement error correlation matrix
$\mathbf{R} = (\rho_{\ell\ell'})$. Consider the sum of centered
measurements:
\begin{equation}
    S_L = \sum_{\ell=1}^{L} (Z_\ell - \E[Z_\ell]).
\end{equation}
Its variance is:
\begin{align}
    \Var(S_L) &= \sum_{\ell=1}^{L} \Var(Z_\ell) + \sum_{\ell \neq \ell'} \Cov(Z_\ell, Z_{\ell'}) \\
              &\leq \sum_{\ell=1}^{L} \frac{1}{4} + \sum_{\ell \neq \ell'} \rho_{\ell\ell'} \sqrt{\Var(Z_\ell) \Var(Z_{\ell'})} \\
              &\leq \frac{L}{4} + \frac{L(L-1)}{4} \bar{\rho}
               = \frac{L}{4}(1 + (L-1)\bar{\rho}).
\end{align}
If the $L$ laboratories were independent with the same total variance,
we would need $L_{\eff}$ laboratories where $L_{\eff}/4 = (L/4)(1 + (L-1)\bar{\rho})$,
yielding $L_{\eff} = L / (1 + (L-1)\bar{\rho})$, as claimed in~\eqref{eq:L_eff}.

\medskip\noindent\textit{Step 3: Construction of effective independent measurements.}
Since $\Var(S_{L}) / \Var(S_{L_{\eff}}^{\text{indep}}) = 1$ by construction,
we can bound the tail probability of the correlated sum by the tail
probability of an independent sum with $L_{\eff}$ terms. Formally,
let $\widetilde{Z}_1, \dots, \widetilde{Z}_{\lfloor L_{\eff}\rfloor}$
be independent random variables with $\E[\widetilde{Z}_i] = 1/2 + \Delta$
and each bounded in $[0,1]$. Then for any $t > 0$:
\begin{equation}
    \Pbb(S_L \leq -t) \leq \Pbb\left(\sum_{i=1}^{\lfloor L_{\eff} \rfloor} (\widetilde{Z}_i - \E[\widetilde{Z}_i]) \leq -t\right).
\end{equation}
This follows from the fact that the variance is an upper bound for tail
probabilities under bounded differences — the correlated sum has the same
variance as $L_{\eff}$ independent terms, making its concentration no
\textit{worse} than the independent case for the same effective count.

\medskip\noindent\textit{Step 4: Hoeffding concentration.}
Define the centered variables $D_i = \widetilde{Z}_i - (1/2 + \Delta)$ with
$|D_i| \leq 1$. The partial sums $S_k = \sum_{i=1}^{k} D_i$ form a martingale
with respect to the natural filtration. By Hoeffding's inequality:
\begin{align}
    \Pbb\left(\sum_{i=1}^{\lfloor L_{\eff}\rfloor} \widetilde{Z}_i \leq \frac{\lfloor L_{\eff}\rfloor}{2} + \lfloor L_{\eff}\rfloor \delta\right)
    &= \Pbb\left(S_{\lfloor L_{\eff}\rfloor} \leq -\lfloor L_{\eff}\rfloor (\Delta - \delta)\right) \\
    &\leq \exp\!\left(-\frac{2 \lfloor L_{\eff}\rfloor^2 (\Delta - \delta)^2}
                           {\sum_{i=1}^{\lfloor L_{\eff}\rfloor} 1^2}\right) \\
    &= \exp\!\bigl(-2 \lfloor L_{\eff}\rfloor (\Delta - \delta)^2\bigr) \\
    &\leq \exp\!\bigl(-2 L_{\eff} (\Delta - \delta)^2\bigr),
\end{align}
for $\Delta > \delta$. The factor of 2 in the theorem statement accounts
for the two-sided nature of the deviation (both overestimation and
underestimation), yielding the bound~\eqref{eq:hoeffding_bound}.

\medskip\noindent\textit{Step 5: Systematic error detection probability.}
A systematic error of magnitude $> \delta B_Q$ is missed by all laboratories
precisely when each laboratory's measurement falls within $\delta B_Q$ of the
spurious value rather than the true value. This requires the consensus
to deviate by at least $\Delta B_Q$ from the true value (since each lab
has margin $\Delta$, the spurious value must attract all measurements).
Setting $\delta = 0$ (the threshold for any non-zero deviation) yields
the worst-case bound $\exp(-2 L_{\eff} \Delta^2)$ as stated in~\eqref{eq:systematic_bound}.
\end{proof}

\begin{corollary}[Required Number of Laboratories for Certified Auditing 认证审计所需实验室数量]
\label{cor:required_L}
To achieve systematic error detection probability at least $1 - \alpha$
with confidence $1 - \gamma$, the effective number of independent laboratories
must satisfy:
\begin{equation}
    L_{\eff} \geq \frac{1}{2\Delta^2} \log\frac{1}{\alpha}.
\end{equation}
For typical laboratory margins $\Delta \approx 0.2$ and target detection
probability $1-\alpha = 0.99$ ($\alpha = 0.01$), this requires
$L_{\eff} \geq 57.6$. With correlated laboratories ($\bar{\rho} = 0.3$),
the nominal $L$ must satisfy $L/(1 + (L-1) \cdot 0.3) \geq 58$,
yielding $L \geq 184$ — indicating that strongly correlated laboratories
provide limited additional certification value.
\end{corollary}

\begin{corollary}[Value of Protocol Diversity 协议多样性的价值]
\label{cor:protocol_diversity}
When laboratories use different measurement protocols, the effective
number of independent measurements increases. Under Assumption~\ref{ass:A4}
(protocol diversity), the correlation between two laboratories measuring
with distinct protocols satisfies $\rho_{\ell\ell'}^{\text{(protocol)}} = \rho_{\ell\ell'}^{\text{(base)}} \cdot \kappa_{pp'}$,
where $\kappa_{pp'} = \int_0^\infty F_p(\omega) F_{p'}(\omega) S(\omega) d\omega / (\text{norm})$
quantifies the overlap of their spectral filter functions. For protocols
with minimal spectral overlap (e.g., Hahn echo vs XY8), $\kappa_{pp'} \ll 1$,
increasing $L_{\eff}$ toward $L$.
\end{corollary}

\begin{remark}[Physical Interpretation for NV Center Auditing]
\label{rem:nv_interpretation}
Theorem~1 establishes that multi-laboratory auditing provides exponentially
strong guarantees against systematic measurement errors. A ``systematic error''
in \NV{} center characterization might arise from:
\begin{itemize}[nosep]
    \item Miscalibrated microwave power leading to systematic overestimation
    of $\RabiFreq$;
    \item Inadequate magnetic shielding causing $\Ttwo$ underestimation
    due to environmental AC magnetic fields;
    \item Incorrect pulse timing causing reduced Hahn echo contrast;
    \item Shared data analysis software with a systematic fitting bias.
\end{itemize}
The bound $\exp(-2 L_{\eff} \Delta^2)$ shows that even modest $L_{\eff}$
(e.g., $L_{\eff}=5$, $\Delta=0.2$) yields detection probability $>0.98$
for systematic errors. The {\em effective} laboratory count is what matters:
ten laboratories sharing the same commercial magnetometer provide little
more assurance than one.
\end{remark}

% ════════════════════════════════════════════════════════════════════════
\section{Theorem 2: Cross-Protocol \Cercis{} Score Convergence 跨协议Cercis评分收敛}
\label{sec:thm2}
% ════════════════════════════════════════════════════════════════════════

\begin{definition}[Coherence Time Reproducibility 相干时间再现性]
\label{def:coherence_repro}
For quality parameter $Q \in \{\Ttwo, \Tone, \RabiFreq\}$, the reproducibility
across laboratories is:
\begin{equation}
    R_Q(\mathbf{s}) = 1 - \frac{\sigma_Q(\mathbf{s})}{\mu_Q(\mathbf{s})},
\end{equation}
where $\mu_Q(\mathbf{s}) = \frac{1}{|\cL_{\text{audit}}|} \sum_{\ell \in \cL_{\text{audit}}} \widehat{Q}^{(\ell, p^*_\ell)}(\mathbf{s})$
is the consensus mean and $\sigma_Q^2(\mathbf{s}) = \frac{1}{|\cL_{\text{audit}}|-1} \sum_{\ell} (\widehat{Q}^{(\ell, p^*_\ell)} - \mu_Q)^2$
is the inter-laboratory variance. $R_Q = 1$ indicates perfect reproducibility;
$R_Q \to 0$ indicates that inter-laboratory variation dominates the signal.
\end{definition}

\begin{definition}[Readout Fidelity Consistency 读取保真度一致性]
\label{def:readout_consistency}
For readout fidelity $\FidReadout$, measured by laboratories using potentially
different readout techniques (single-shot, time-averaged, spin-to-charge
conversion), the consistency is:
\begin{equation}
    C_{\FidReadout}(\mathbf{s}) = 1 - \max_{\ell, \ell' \in \cL_{\text{audit}}} |\widehat{\FidReadout}^{(\ell)}(\mathbf{s}) - \widehat{\FidReadout}^{(\ell')}(\mathbf{s})|.
\end{equation}
This is the complement of the maximum pairwise discrepancy, a conservative
metric that penalizes any single outlier laboratory.
\end{definition}

\begin{definition}[Fabrication Method Novelty 制备方法新颖性]
\label{def:fab_novelty}
For state $\mathbf{s}$ with fabrication parameters, the novelty score is:
\begin{equation}
    N(\mathbf{s}) = \sum_{k \in \text{fab}} v_k \cdot \ind{s_k \notin \cS_{\text{cal}, k}},
    \label{eq:novelty}
\end{equation}
where $\cS_{\text{cal}, k}$ is the set of parameter values for dimension $k$
present in the calibration set, and $v_k \geq 0$ are dimension weights
with $\sum_k v_k = 1$. $N(\mathbf{s}) = 0$ means all fabrication parameters
are represented in calibration; $N(\mathbf{s}) = 1$ means no calibration
precedent exists. Fabrication dimensions include substrate type, nitrogen
isotope, implantation method, surface treatment, and annealing protocol.
\end{definition}

\begin{definition}[Cercis Score for NV Centers \Cercis 评分]
\label{def:cercis}
For an \NV{} center in state $\mathbf{s}$, the \textbf{\Cercis{} score} is:
\begin{equation}
    S(\mathbf{s}) = Q(\mathbf{s}) + \eta \cdot N(\mathbf{s}),
    \label{eq:cercis_def}
\end{equation}
where:
\begin{itemize}[nosep]
    \item $Q(\mathbf{s}) = \alpha_T \cdot R_{\Ttwo}(\mathbf{s}) + \alpha_F \cdot C_{\FidReadout}(\mathbf{s})$,
    with $\alpha_T + \alpha_F = 1$, $\alpha_T, \alpha_F \geq 0$, is the
    \textbf{quality evidence concordance} (品质证据一致性);
    \item $N(\mathbf{s})$ is the fabrication method novelty from
    Definition~\ref{def:fab_novelty};
    \item $\eta \geq 0$ is the \textbf{novelty weight}.
\end{itemize}
\end{definition}

\begin{theorem}[Cross-Protocol \Cercis{} Score Convergence 跨协议Cercis评分收敛]
\label{thm:cercis_convergence}
\rigorFull

Let $\cL_{\text{audit}} \subset \cL$ be a set of $L_{\text{audit}}$ laboratories
that audit state $\mathbf{s}$ using protocols from $\cP$, satisfying
Assumptions~\ref{ass:A1}--\ref{ass:A8}.

Define the true quality score:
\begin{equation}
    S^*(\mathbf{s}) = \alpha_T \cdot \left(1 - \frac{\sigma_Q^{\text{true}}(\mathbf{s})}{Q(\mathbf{s})}\right)
                   + \alpha_F \cdot \left(1 - \delta_{\FidReadout}^{\text{true}}(\mathbf{s})\right)
                   + \eta \cdot N(\mathbf{s}),
\end{equation}
where $\sigma_Q^{\text{true}}$ is the protocol-limited measurement uncertainty
and $\delta_{\FidReadout}^{\text{true}}$ is the quantum-limited readout infidelity.

Then the empirical \Cercis{} score $\widehat{S}(\mathbf{s})$ computed from
$L_{\text{audit}}$ laboratories converges to $S^*(\mathbf{s})$ with rate:
\begin{equation}
    \Pbb\left(|\widehat{S}(\mathbf{s}) - S^*(\mathbf{s})| > \varepsilon\right)
    \leq 4 \exp\!\left(-\frac{L_{\eff} \cdot \varepsilon^2}{8}\right)
        + \cO\!\left(\frac{1}{\sqrt{n_{\text{cal}}}}\right),
    \label{eq:cercis_convergence}
\end{equation}
where $L_{\eff}$ is defined as in Theorem~\ref{thm:error_detection} and
$n_{\text{cal}}$ is the calibration set size from Assumption~\ref{ass:A5}.
\end{theorem}

\begin{proof}
We decompose the estimation error into laboratory sampling error and
calibration estimation error.

\medskip\noindent\textit{Step 1: Decomposition of the estimation error.}
The empirical \Cercis{} score has three components:
\begin{equation}
    \widehat{S}(\mathbf{s}) = \alpha_T \widehat{R}_{\Ttwo}(\mathbf{s}) + \alpha_F \widehat{C}_{\FidReadout}(\mathbf{s}) + \eta N(\mathbf{s}).
\end{equation}
The novelty term $N(\mathbf{s})$ is deterministic given $\mathbf{s}$ (it
depends only on whether fabrication parameters are in the calibration
set), so the only error sources are in $\widehat{R}_{\Ttwo}$ and $\widehat{C}_{\FidReadout}$.

The estimation error decomposes as:
\begin{equation}
    |\widehat{S} - S^*| \leq \alpha_T |\widehat{R}_{\Ttwo} - R_{\Ttwo}^*| + \alpha_F |\widehat{C}_{\FidReadout} - C_{\FidReadout}^*|,
\end{equation}
with $R_{\Ttwo}^* = 1 - \sigma_{\Ttwo}^{\text{true}} / \Ttwo(\mathbf{s})$ and
$C_{\FidReadout}^* = 1 - \delta_{\FidReadout}^{\text{true}}$.

\medskip\noindent\textit{Step 2: Concentration of $\widehat{R}_{\Ttwo}$.}
The inter-laboratory variance estimator is:
\begin{equation}
    \widehat{\sigma}_{\Ttwo}^2 = \frac{1}{L_{\text{audit}}-1} \sum_{\ell=1}^{L_{\text{audit}}} (\widehat{T}_2^{(\ell)} - \widehat{\mu}_{\Ttwo})^2.
\end{equation}
By the bounded differences property (Assumption~\ref{ass:A1}) and Hoeffding's
inequality for U-statistics, the variance estimator concentrates:
\begin{equation}
    \Pbb(|\widehat{\sigma}_{\Ttwo}^2 - (\sigma_{\Ttwo}^{\text{true}})^2| > \varepsilon_1)
    \leq 2\exp\!\left(-\frac{L_{\eff} \cdot \varepsilon_1^2}{2 B_{\Ttwo}^4}\right).
\end{equation}
Similarly, the mean estimator concentrates:
\begin{equation}
    \Pbb(|\widehat{\mu}_{\Ttwo} - \Ttwo(\mathbf{s})| > \varepsilon_2)
    \leq 2\exp\!\left(-\frac{L_{\eff} \cdot \varepsilon_2^2}{2 B_{\Ttwo}^2}\right).
\end{equation}

\medskip\noindent\textit{Step 3: Propagation to $\widehat{R}_{\Ttwo}$.}
By the delta method, since $R = 1 - \sigma / \mu$ with both $\sigma$ and $\mu$
concentrating, the coefficient of variation (inverse of $R$) concentrates
with the same exponential rate. Setting $\varepsilon_1 = \varepsilon_2 = \varepsilon / 4$
and applying the union bound over the two concentration events yields:
\begin{equation}
    \Pbb(|\widehat{R}_{\Ttwo} - R_{\Ttwo}^*| > \varepsilon) \leq 2\exp\!\left(-\frac{L_{\eff} \cdot \varepsilon^2}{8 \max(B_{\Ttwo}^2, B_{\Ttwo}^4)}\right).
\end{equation}
For $B_{\Ttwo} \geq 1$ (the physically relevant regime where coherence times
exceed measurement precision by at least one order of magnitude), the
denominator simplifies to $8 B_{\Ttwo}^4$; normalizing by setting $B_{\Ttwo}=1$
via rescaling yields the form in~\eqref{eq:cercis_convergence}.

\medskip\noindent\textit{Step 4: Concentration of $\widehat{C}_{\FidReadout}$.}
The maximum pairwise discrepancy is a $\cO(1/\sqrt{L_{\text{audit}}})$ estimator
of the population range. By the bounded differences inequality applied to
the readout fidelity measurements (each in $[0,1]$):
\begin{equation}
    \Pbb(|\widehat{C}_{\FidReadout} - C_{\FidReadout}^*| > \varepsilon)
    \leq 2\exp\!\left(-2 L_{\eff} \varepsilon^2\right).
\end{equation}

\medskip\noindent\textit{Step 5: Calibration error.}
The weights $\alpha_T, \alpha_F$ and novelty weights $v_k$ are estimated
from the calibration set $\cS_{\text{cal}}$. By standard empirical process
theory, their estimation error is $\cO(1/\sqrt{n_{\text{cal}}})$, which
adds a negligible term for sufficiently large $n_{\text{cal}}$.

\medskip\noindent\textit{Step 6: Union bound.}
Combining Steps 2--5 via the union bound (each of the 4 tail events for
$\sigma$, $\mu$, $C_{\FidReadout}$, and calibration), we obtain:
\begin{equation}
    \Pbb(|\widehat{S} - S^*| > \varepsilon) \leq 4\exp\!\left(-\frac{L_{\eff} \varepsilon^2}{8}\right) + \cO(n_{\text{cal}}^{-1/2}),
\end{equation}
as stated.
\end{proof}

\begin{corollary}[Cercis Score as Audit Quality Certificate]
\label{cor:cercis_cert}
For a target certification confidence $1 - \gamma$, the \Cercis{} score
$\widehat{S}(\mathbf{s})$ certifies that the true quality $S^*(\mathbf{s})$
lies within $\widehat{S}(\mathbf{s}) \pm \varepsilon$ with:
\begin{equation}
    \varepsilon = \sqrt{\frac{8}{L_{\eff}} \log\frac{4}{\gamma}} + \cO(n_{\text{cal}}^{-1/2}).
\end{equation}
For $L_{\eff}=10$ and $\gamma = 0.05$, this yields $\varepsilon \approx 1.35$
in absolute score units — indicating that certification at high precision
requires larger effective laboratory counts or protocol diversity.
\end{corollary}

\begin{remark}
Theorem~2 establishes that the \Cercis{} score is not merely a heuristic
but a statistically consistent estimator of a well-defined true quality
metric. The convergence rate depends on $L_{\eff}$ (not merely $L$),
emphasizing that correlated laboratories provide diminishing certification
returns — a finding with direct policy implications for \NV{} center
quality standardization efforts.
\end{remark}

% ════════════════════════════════════════════════════════════════════════
\section{Theorem 3: $\Ttwo$ Degradation Source Unidentifiability $\Ttwo$退化源不可辨识性}
\label{sec:thm3}
% ════════════════════════════════════════════════════════════════════════

\begin{theorem}[$\Ttwo$ Degradation Source Unidentifiability 相干时间退化源不可辨识性]
\label{thm:unident}
\rigorFull

Let two laboratories $\ell_A, \ell_B$ measure the coherence time of \NV{}
center samples in states $\mathbf{s}_A$ and $\mathbf{s}_B$, respectively.
Suppose the laboratories report $\widehat{T}_2^{(\ell_A)}(\mathbf{s}_A) \neq \widehat{T}_2^{(\ell_B)}(\mathbf{s}_B)$,
i.e., a discrepancy in measured coherence times.

Consider three possible causal sources of the discrepancy:
\begin{enumerate}[label=(E\arabic*), nosep]
    \item $E_1$: \textbf{Surface noise dominance} — one sample has higher
    surface spin noise due to different surface termination or shallower
    \NV{} implantation depth, causing genuine $\Ttwo$ degradation;
    \item $E_2$: \textbf{Bulk impurity dominance} — one sample has higher
    bulk paramagnetic impurity concentration (P1 centers, ${}^{13}$C nuclear
    spins), causing genuine $\Ttwo$ degradation;
    \item $E_3$: \textbf{Measurement protocol artifact} — the discrepancy
    arises not from sample differences but from the laboratories using
    different measurement protocols with distinct spectral filter functions,
    sensitivity to environmental noise, or calibration procedures.
\end{enumerate}

Then, without at least one of the following declared structural assumptions —
\begin{enumerate}[label=(S\arabic*), nosep]
    \item Known depth profile of \NV{} centers in both samples, including
    implantation energy, annealing conditions, and surface characterization
    (XPS, SIMS);
    \item Known bulk impurity concentrations (EPR for P1 centers, SIMS for
    nitrogen, Raman for ${}^{13}$C isotopic abundance) for both samples;
    \item Declared measurement protocol specifications: pulse sequence,
    filter function $F_p(\omega)$, microwave power calibration, laser
    repumping scheme, and environmental noise spectrum $S(\omega)$ at
    each laboratory —
\end{enumerate}
the true causal source among $\{E_1, E_2, E_3\}$ is \textbf{fundamentally
unidentifiable} from the observed $\Ttwo$ discrepancy alone.

Furthermore, for any desired posterior distribution $\mathbf{q} = (q_1, q_2, q_3)$
over the three sources (with $q_i \geq 0$, $\sum_i q_i = 1$), there exists
a consistent pair of measurement configurations $(\ell_A, \mathbf{s}_A, p_A)$
and $(\ell_B, \mathbf{s}_B, p_B)$ that produces the observed $\Ttwo$
discrepancy while having exactly that posterior over causal sources.
\end{theorem}

\begin{proof}
We construct three observationally equivalent worlds, each attributing
the same $\Ttwo$ discrepancy to a different dominant cause. This is a
proof by construction of non-identifiability in the latent variable model.

\medskip\noindent\textit{Step 1: Latent variable specification.}
Define three binary indicator variables:
\begin{align}
    Z_1 &= \ind{\Ttwo \text{ discrepancy dominated by surface noise } (E_1)}, \\
    Z_2 &= \ind{\Ttwo \text{ discrepancy dominated by bulk impurities } (E_2)}, \\
    Z_3 &= \ind{\Ttwo \text{ discrepancy dominated by protocol artifact } (E_3)}.
\end{align}
These are mutually exclusive: $Z_1 + Z_2 + Z_3 = 1$.

\medskip\noindent\textit{Step 2: Observation model.}
The observed data is the ordered pair of measured coherence times:
\begin{equation}
    \mathbf{o} = (\widehat{T}_2^{(\ell_A)}(\mathbf{s}_A), \widehat{T}_2^{(\ell_B)}(\mathbf{s}_B)) \in \R_+^2,
\end{equation}
with $\widehat{T}_2^{(\ell_A)} \neq \widehat{T}_2^{(\ell_B)}$. Define the
discrepancy magnitude $\delta = |\widehat{T}_2^{(\ell_A)} - \widehat{T}_2^{(\ell_B)}| > 0$.

The mapping from latent causes to observed measurements is parameterized by:
\begin{align}
    \theta_i &= (\Ttwo(\mathbf{s}_A; Z_i=1), \Ttwo(\mathbf{s}_B; Z_i=1),
                 \beta_{\ell_A}(p_A, \mathbf{s}_A; Z_i=1),
                 \beta_{\ell_B}(p_B, \mathbf{s}_B; Z_i=1)), \quad i \in \{1,2,3\},
\end{align}
where $\beta$ are systematic biases. The observation likelihood given
cause $Z_i = 1$ is:
\begin{equation}
    \mathcal{L}_i = \prod_{m \in \{A,B\}} \phi\!\left(\frac{\widehat{T}_2^{(\ell_m)} - \Ttwo(\mathbf{s}_m; Z_i=1) - \beta_{\ell_m}}{\sigma_{\ell_m}}\right),
\end{equation}
where $\phi$ is the standard normal density.

\medskip\noindent\textit{Step 3: World 1 — Surface noise dominates ($Z_1 = 1$).}
Choose $\mathbf{s}_A$ to be a shallow \NV{} center (depth $d_A = 5$~nm) with
hydrogen-terminated surface, and $\mathbf{s}_B$ to be a deep \NV{} center
(depth $d_B = 50$~nm) in the same diamond substrate. The surface noise
spectral density $S_{\text{surf}}(\omega) \propto 1/d^2$ makes $\Ttwo(\mathbf{s}_A) \ll \Ttwo(\mathbf{s}_B)$.
Both laboratories use Hahn echo ($p_A = p_B = \text{Hahn}$) with identical
apparatus ($\beta_A = \beta_B = 0$). The observed discrepancy $\delta$ equals
$\Ttwo(\mathbf{s}_B) - \Ttwo(\mathbf{s}_A)$, purely from surface noise.

With parameter choices $\Ttwo(\mathbf{s}_A; Z_1{=}1) = 100~\mu\text{s}$,
$\Ttwo(\mathbf{s}_B; Z_1{=}1) = 500~\mu\text{s}$, identical protocols
and calibration, the likelihood is:
\begin{equation}
    \mathcal{L}_1 = \phi(0) \cdot \phi(0) = \frac{1}{2\pi\sigma^2}.
\end{equation}

\medskip\noindent\textit{Step 4: World 2 — Bulk impurities dominate ($Z_2 = 1$).}
Choose $\mathbf{s}_A$ to be an \NV{} center in HPHT diamond with
$[\text{N}] = 100$~ppm, and $\mathbf{s}_B$ to be an \NV{} center in
electronic-grade CVD diamond with $[\text{N}] < 1$~ppb. Both centers are
at the same depth ($d_A = d_B = 50$~nm, eliminating surface noise),
both laboratories use Hahn echo. The paramagnetic spin bath density in
$\mathbf{s}_A$ causes $\Ttwo(\mathbf{s}_A) \ll \Ttwo(\mathbf{s}_B)$.

With parameter choices $\Ttwo(\mathbf{s}_A; Z_2{=}1) = 100~\mu\text{s}$
(identical value as World~1), $\Ttwo(\mathbf{s}_B; Z_2{=}1) = 500~\mu\text{s}$
(identical value as World~1), the likelihood is:
\begin{equation}
    \mathcal{L}_2 = \phi(0) \cdot \phi(0) = \frac{1}{2\pi\sigma^2} = \mathcal{L}_1.
\end{equation}

\medskip\noindent\textit{Step 5: World 3 — Protocol artifact dominates ($Z_3 = 1$).}
Choose $\mathbf{s}_A = \mathbf{s}_B$ (identical \NV{} center, eliminating
both surface and bulk heterogeneity). Let laboratory A use Hahn echo while
laboratory B uses XY8 dynamical decoupling. The true $\Ttwo$ is identical
for both samples, but the different filter functions $F_{\text{Hahn}}(\omega)$
and $F_{\text{XY8}}(\omega)$ produce different measured values:
$\widehat{T}_2^{(\ell_A)} = \Ttwo_{\text{true}} + \delta_{\text{Hahn}}$,
$\widehat{T}_2^{(\ell_B)} = \Ttwo_{\text{true}} + \delta_{\text{XY8}}$,
with $\delta_{\text{XY8}} > \delta_{\text{Hahn}}$.

With parameter choices $\Ttwo(\mathbf{s}_A; Z_3{=}1) = \Ttwo(\mathbf{s}_B; Z_3{=}1) = 300~\mu\text{s}$,
$\beta_A = -200~\mu\text{s}$ (Hahn echo systematic underestimation relative
to true $\Ttwo$), $\beta_B = +200~\mu\text{s}$ (XY8 systematic overestimation),
the observed measurements are $\widehat{T}_2^{(\ell_A)} = 100~\mu\text{s}$
and $\widehat{T}_2^{(\ell_B)} = 500~\mu\text{s}$ — identical to Worlds~1 and~2.
The likelihood is:
\begin{equation}
    \mathcal{L}_3 = \phi(0) \cdot \phi(0) = \frac{1}{2\pi\sigma^2} = \mathcal{L}_1 = \mathcal{L}_2.
\end{equation}

\medskip\noindent\textit{Step 6: General impossibility of causal attribution.}
The three worlds produce identical observations $\mathbf{o} = (100~\mu\text{s}, 500~\mu\text{s})$
with identical likelihoods $\mathcal{L}_1 = \mathcal{L}_2 = \mathcal{L}_3$,
yet each attributes the discrepancy to a fundamentally different physical cause.
The parameter space has 12 free parameters (4 parameters per cause $\times$ 3 causes),
constrained by one observed discrepancy $\delta$. The solution manifold has
dimension $12 - 1 - 2 = 9$ (accounting for the measurement pair and the
probability simplex constraints), meaning there is a 9-dimensional family of
configurations that produce exactly the same observations.

For any desired posterior $\mathbf{q}^*$, we can interpolate between the
three extremal worlds to achieve it: set $\Ttwo(\mathbf{s}_A) = q_1^* \cdot 100 + q_2^* \cdot 100 + q_3^* \cdot 300$,
$\Ttwo(\mathbf{s}_B) = q_1^* \cdot 500 + q_2^* \cdot 500 + q_3^* \cdot 300$,
and similarly for biases, preserving the observed measurements while
realizing the specified posterior.

\medskip\noindent\textit{Step 7: Necessary conditions for identifiability.}
Each structural assumption (S1)--(S3) breaks the observational equivalence:
\begin{itemize}[nosep]
    \item (S1): Known depth profiles distinguish World~1 (different depths)
    from Worlds~2 and~3 (same depth);
    \item (S2): Known impurity concentrations distinguish World~2 (different
    [N]) from Worlds~1 and~3;
    \item (S3): Declared protocols distinguish World~3 (different protocols)
    from Worlds~1 and~2 (same protocol).
\end{itemize}
Without at least one, the three worlds remain indistinguishable.
Without all three, at best pairs of causes can be conflated.
\end{proof}

\begin{corollary}[Practical Implication for NV Center Publication Standards]
\label{cor:publication_standard}
Theorem~3 implies that any publication reporting an \NV{} center $\Ttwo$
value must, at minimum, declare:
\begin{enumerate}[label=(\roman*), nosep]
    \item The \NV{} implantation depth (or a measured depth profile);
    \item The bulk nitrogen and ${}^{13}$C concentrations of the diamond substrate;
    \item The complete measurement protocol specification, including pulse
    sequence, repetition rate, laser power, and environmental noise characterization.
\end{enumerate}
Without all three declarations, the reported $\Ttwo$ value cannot be
causally interpreted — it could reflect genuine material quality, or
it could be an artifact of measurement choices. This is a minimum
information standard (\textbf{NV-MIS}) for \NV{} center characterization.
\end{corollary}

\begin{remark}[Relation to Derman's Principle of Model Risk]
Theorem~3 embodies, in the quantum materials context, Derman's framework
of model risk~\cite{derman1996model}: the source of a discrepancy between
measurement and expectation (here, between two laboratories' $\Ttwo$
measurements) cannot be resolved without declared assumptions about the
generative process. Our construction of three observationally equivalent
worlds formalizes this principle for \NV{} center auditing.
\end{remark}

% ════════════════════════════════════════════════════════════════════════
\section{\Cercis{} Score and \Yajie{} Multi-Laboratory Consensus \Cercis 评分与Yajie多实验室共识}
\label{sec:cercis}
% ════════════════════════════════════════════════════════════════════════

\subsection{\Cercis{} Score: Detailed Formulation}

The \Cercis{} score $S(\mathbf{s}) = Q(\mathbf{s}) + \eta N(\mathbf{s})$ was
defined in Definition~\ref{def:cercis}. We now provide the fully operationalized
form for \NV{} center quality auditing.

\begin{definition}[Operationalized Quality Concordance 可操作的品质一致性]
\label{def:Q_operational}
For state $\mathbf{s}$ audited by laboratories $\cL_{\text{audit}}$:
\begin{equation}
    Q(\mathbf{s}) = \alpha_T \cdot R_{\Ttwo}(\mathbf{s})
                 + \alpha_F \cdot C_{\FidReadout}(\mathbf{s})
                 + \alpha_\Omega \cdot R_{\RabiFreq}(\mathbf{s}),
    \label{eq:Q_full}
\end{equation}
with $\alpha_T + \alpha_F + \alpha_\Omega = 1$, where:
\begin{itemize}[nosep]
    \item $R_{\Ttwo}(\mathbf{s})$ is the $\Ttwo$ reproducibility
    (Definition~\ref{def:coherence_repro});
    \item $C_{\FidReadout}(\mathbf{s})$ is the readout fidelity consistency
    (Definition~\ref{def:readout_consistency});
    \item $R_{\RabiFreq}(\mathbf{s})$ is the Rabi frequency reproducibility,
    defined analogously to $R_{\Ttwo}$;
    \item Default weights: $\alpha_T = 0.5, \alpha_F = 0.3, \alpha_\Omega = 0.2$,
    reflecting the primacy of coherence time as the most information-rich
    \NV{} quality parameter.
\end{itemize}
\end{definition}

\begin{definition}[Operationalized Fabrication Novelty 可操作的制备新颖性]
\label{def:N_operational}
The fabrication novelty $N(\mathbf{s})$ decomposes across six dimensions,
each weighted by its impact on \NV{} center performance:
\begin{equation}
    N(\mathbf{s}) = \sum_{k=1}^{6} v_k \cdot n_k(\mathbf{s}),
    \label{eq:N_full}
\end{equation}
where:
\begin{center}
\begin{tabular}{lll}
\toprule
\textbf{Dimension} & \textbf{Weight $v_k$} & \textbf{Novelty indicator $n_k(\mathbf{s})$} \\
\midrule
Substrate type (CVD/HPHT) & 0.20 & $\ind{s_{\text{sub}} \notin \cS_{\text{cal}, \text{sub}}}$ \\
Nitrogen isotope (N-14/N-15) & 0.15 & $\ind{s_{\text{iso}} \notin \cS_{\text{cal}, \text{iso}}}$ \\
Implantation depth regime & 0.25 & $\ind{s_{\text{depth}} \notin \cS_{\text{cal}, \text{depth}}}$ \\
Surface termination chemistry & 0.15 & $\ind{s_{\text{surf}} \notin \cS_{\text{cal}, \text{surf}}}$ \\
${}^{13}$C isotopic enrichment & 0.15 & $\ind{s_{\text{C13}} \notin \cS_{\text{cal}, \text{C13}}}$ \\
Annealing/fabrication protocol & 0.10 & $\ind{\text{protocol} \notin \cS_{\text{cal}, \text{fab}}}$ \\
\bottomrule
\end{tabular}
\end{center}
\end{definition}

\begin{proposition}[Monotonicity of \Cercis{} Score]
\label{prop:cercis_monotone}
Under Assumptions~\ref{ass:A1}--\ref{ass:A8}, the \Cercis{} score satisfies:
\begin{enumerate}[label=(\roman*), nosep]
    \item $S(\mathbf{s}) \in [0, 1 + \eta]$ for all $\mathbf{s} \in \cS$;
    \item $S$ is non-decreasing in each laboratory's measurement precision;
    \item $S(\mathbf{s}) \to 1 + \eta N(\mathbf{s})$ as $L_{\eff} \to \infty$
    and all laboratories achieve their asymptotic measurement limits.
\end{enumerate}
\end{proposition}

\begin{proof}
(i) Follows from $R_{\Ttwo}, C_{\FidReadout}, R_{\RabiFreq} \in [0,1]$
and the convex combination weights summing to 1, with $N(\mathbf{s}) \in [0,1]$.

(ii) Measurement precision enters through $\sigma_{\ell, p}$ in the
measurement model~\eqref{eq:measurement_model}. As precision improves
($\sigma_{\ell, p} \to 0$), the inter-laboratory variance decreases,
increasing $R_Q$ monotonically.

(iii) As $L_{\eff} \to \infty$, Theorem~2 implies $\widehat{S} \to S^*$,
and with perfect precision $S^* = 1 + \eta N(\mathbf{s})$ (since
$\sigma_Q^{\text{true}} \to 0$ and $\delta_{\FidReadout}^{\text{true}} \to 0$).
\end{proof}

\subsection{\Yajie{} Multi-Laboratory Consensus}

\begin{definition}[\Yajie{} Multi-Laboratory Consensus 多实验室共识]
\label{def:yajie}
Given $L$ laboratories with measurements $\{\widehat{Q}^{(\ell, p^*_\ell)}(\mathbf{s})\}_{\ell=1}^{L}$
for quality parameter $Q$ on state $\mathbf{s}$, the \textbf{\Yajie{} consensus}
estimate is:
\begin{equation}
    \widehat{Q}_{\Yajie}(\mathbf{s}) = \sum_{\ell=1}^{L} w_\ell(\mathbf{s}) \cdot \widehat{Q}^{(\ell, p^*_\ell)}(\mathbf{s}),
    \label{eq:yajie_consensus}
\end{equation}
where the weights are:
\begin{equation}
    w_\ell(\mathbf{s}) = \frac{S(\mathbf{s}) \cdot \omega_\ell + (1 - S(\mathbf{s})) \cdot \bar{\omega}}
                          {\sum_{\ell'=1}^{L} \bigl(S(\mathbf{s}) \cdot \omega_{\ell'} + (1 - S(\mathbf{s})) \cdot \bar{\omega}\bigr)},
    \label{eq:yajie_weights}
\end{equation}
with:
\begin{itemize}[nosep]
    \item $\omega_\ell$: base weight for laboratory $\ell$, computed as
    the inverse of its mean squared error on the calibration set $\cS_{\text{cal}}$:
    $\omega_\ell = \left(\sum_{\mathbf{s} \in \cS_{\text{cal}}} (\widehat{Q}^{(\ell, p^*_\ell)}(\mathbf{s}) - Q_{\text{ref}}(\mathbf{s}))^2\right)^{-1}$;
    \item $\bar{\omega} = \frac{1}{L} \sum_{\ell=1}^{L} \omega_\ell$: uniform baseline weight.
\end{itemize}

\textbf{Weighting principle}: When the \Cercis{} score $S(\mathbf{s})$ is high
(high reproducibility, low novelty — the sample is ``well-understood''),
weights concentrate on precise laboratories. When $S(\mathbf{s})$ is low
(poor reproducibility or high novelty — the sample is ``unfamiliar''),
weights revert to uniform, preventing overconfident extrapolation.
\end{definition}

\begin{proposition}[\Yajie{} Consensus Optimality Property]
\label{prop:yajie_optimal}
Under Assumptions~\ref{ass:A1}--\ref{ass:A5}, the \Yajie{} consensus estimator
minimizes the expected squared error among all linear estimators with
weights in the convex hull of $\{\omega_\ell\}$ and $\bar{\omega}$:
\begin{equation}
    \E[(\widehat{Q}_{\Yajie} - Q(\mathbf{s}))^2] \leq \min_{\mathbf{w} \in \text{conv}(\boldsymbol{\omega}, \bar{\omega}\mathbf{1})} \E[(\widehat{Q}_{\mathbf{w}} - Q(\mathbf{s}))^2] + \cO(L_{\eff}^{-1}).
\end{equation}
\end{proposition}

\begin{proof}
The \Yajie{} weights~\eqref{eq:yajie_weights} are a convex combination of
the precision-optimal weights $\boldsymbol{\omega}$ (minimizing variance
when biases are zero) and the uniform weights $\bar{\omega}\mathbf{1}$
(minimizing worst-case bias when biases are unknown). By the fundamental
bias-variance decomposition, this achieves near-optimal risk for the minimax
problem over the uncertainty class defined by $S(\mathbf{s})$, with excess
risk $\cO(L_{\eff}^{-1})$ from finite-sample weight estimation.
\end{proof}

% ════════════════════════════════════════════════════════════════════════
\section{Experimental Benchmarks 实验基准}
\label{sec:experiments}
% ════════════════════════════════════════════════════════════════════════

The experimental framework evaluates the \SCX{} auditing theorems across
four benchmark axes, each probing a distinct dimension of \NV{} center
heterogeneity.

\subsection{Benchmark Design}

\begin{enumerate}[label=\textbf{B\arabic*:}, leftmargin=*]

    \item \textbf{N-14 vs N-15 Isotopic Composition (N-14与N-15同位素对比).}
    \begin{itemize}[nosep]
        \item \textit{Motivation}: N-14 has nuclear spin $I=1$ (three-level
        hyperfine structure), while N-15 has $I=1/2$ (two-level). The hyperfine
        coupling strength differs ($A_{\text{N-14}} \approx 2.16$~MHz vs
        $A_{\text{N-15}} \approx 3.03$~MHz), affecting decoherence channels
        and readout fidelity.
        \item \textit{Protocol}: Prepare matched pairs of \NV{} centers — one
        ensemble with natural N-14, one with isotopically enriched N-15 — in
        identical CVD diamond substrates at matching depths. Distribute to
        $L \geq 5$ laboratories for blind measurement of $\Ttwo$, $\TtwoStar$,
        $\RabiFreq$, and $\FidReadout$.
        \item \textit{SCX Metrics}: Compute $R_{\Ttwo}$ for N-14 and N-15
        cohorts separately. Test whether the \Yajie{} consensus detects
        systematic differences in $\Ttwo$ that are attributable to isotopic
        composition rather than measurement noise (Theorem~1 bound application).
    \end{itemize}

    \item \textbf{Shallow vs Deep NV Centers (浅层与深层NV中心对比).}
    \begin{itemize}[nosep]
        \item \textit{Motivation}: Shallow \NV{} centers ($d < 10$~nm) are
        essential for nanoscale NMR sensing but suffer from surface-induced
        decoherence. Deep \NV{} centers ($d > 50$~nm) are bulk-noise-limited.
        The surface-to-bulk noise transition tests Theorem~3 unidentifiability.
        \item \textit{Protocol}: Implant \NV{} centers at 5 depths ($d \in \{3, 8, 15, 30, 80\}$~nm)
        in the same CVD diamond substrate with identical surface termination.
        Each laboratory measures all depths using Hahn echo. Additionally,
        each laboratory characterizes surface noise via $\Tone$ relaxometry
        at variable depth (structural assumption S1 satisfaction).
        \item \textit{SCX Metrics}: Fit $1/\Ttwo(d) = 1/\Ttwo^{\text{bulk}} + \gamma/d^2$
        (surface noise model). Measure the inter-laboratory variance of the
        fitted $\Ttwo^{\text{bulk}}$ parameter. Low variance certifies that
        surface noise is correctly accounted for (validating S1).
    \end{itemize}

    \item \textbf{CVD vs HPHT Diamond Substrates (CVD与HPHT金刚石衬底对比).}
    \begin{itemize}[nosep]
        \item \textit{Motivation}: CVD and HPHT diamond differ systematically
        in nitrogen incorporation, strain distribution, and defect density.
        CVD diamond can achieve $[\text{N}] < 1$~ppb, while HPHT diamond
        typically has $[\text{N}] \sim 1$--$200$~ppm. These differences
        directly impact $\Ttwo$.
        \item \textit{Protocol}: Procure matched pairs: (i) electronic-grade
        CVD ($[\text{N}] < 5$~ppb, ${}^{13}$C $< 0.01\%$), (ii) standard
        CVD ($[\text{N}] \sim 10$--$50$~ppb), (iii) type-IIa HPHT
        ($[\text{N}] < 1$~ppm), (iv) type-Ib HPHT ($[\text{N}] \sim 100$~ppm).
        All with deep \NV{} centers ($d = 50$~nm) to eliminate surface effects.
        Audit by $L \geq 5$ laboratories.
        \item \textit{SCX Metrics}: \Cercis{} score $S(\mathbf{s})$ for each
        substrate type. $Q$ component should be highest for electronic-grade
        CVD (best reproducibility), while $N$ component captures the novelty
        of isotopically engineered substrates.
    \end{itemize}

    \item \textbf{Cross-Protocol Measurement Comparison (跨协议测量对比).}
    \begin{itemize}[nosep]
        \item \textit{Motivation}: Different measurement protocols probe
        different noise spectral regions (Assumption~\ref{ass:A4}). The
        relationship between $\TtwoStar$ (ODMR), $\Ttwo$ (Hahn echo), and
        $T_{2,\text{DD}}$ (XY8) provides information about the noise spectrum
        $S(\omega)$.
        \item \textit{Protocol}: For a single \NV{} center sample (state fixed),
        each laboratory performs all four protocols: ODMR ($\TtwoStar$),
        Ramsey ($\TtwoStar$ with different filter), Hahn echo ($\Ttwo$),
        and XY8 ($T_{2,\text{XY8}}$). Report all four values with measurement
        uncertainties.
        \item \textit{SCX Metrics}: The protocol amplification ratio
        $\mathcal{A} = T_{2,\text{XY8}} / \TtwoStar$ quantifies the noise
        spectrum slope. Inter-laboratory consistency of $\mathcal{A}$ is a
        protocol-independent quality metric. If $\mathcal{A}$ varies across
        laboratories on the same sample, this indicates unaccounted systematic
        biases (Theorem~1 detection). The unidentifiability of whether
        $\mathcal{A}$ variation arises from protocol implementation differences
        or genuine environmental noise differences is established by Theorem~3.
    \end{itemize}

\end{enumerate}

\subsection{Evaluation Metrics}

\begin{enumerate}[label=\textbf{M\arabic*:}, leftmargin=*]
    \item \textbf{Consensus Error Rate.} Compare $\Pbb(|\widehat{Q}_{\Yajie} - Q_{\text{ref}}| > \delta)$
    against the Theorem~1 bound $\exp(-2 L_{\eff} \Delta^2)$. For calibration
    samples with known reference values $Q_{\text{ref}}$, measure the empirical
    tail probability and test whether it respects the Hoeffding bound.

    \item \textbf{Effective Laboratory Count $L_{\eff}$.} Estimate
    $\bar{\rho}$ from the empirical error correlation matrix on calibration
    samples and compute $L_{\eff} = L / (1 + (L-1)\bar{\rho})$. Measure how
    $L_{\eff}$ degrades when laboratories share equipment vendors, calibration
    standards, or data analysis software.

    \item \textbf{\Cercis{} Score Calibration.} For samples with known
    fabrication parameters and independently verified quality (e.g., by
    multi-lab consensus on calibration set), compute $\Gamma(S) = \E[|\E[Q_{\text{ref}} \mid S] - S|]$.
    A well-calibrated \Cercis{} score has $\Gamma(S) \approx 0$.

    \item \textbf{Unidentifiability Verification.} For the depth-series
    benchmark (B2), test whether laboratories can distinguish surface-dominated
    from bulk-dominated decoherence using only $\Ttwo$ measurements (without
    S1 depth information). The prediction from Theorem~3 is that they cannot.

    \item \textbf{Novelty Detection.} For diamond samples fabricated by
    novel methods (new surface treatments, exotic substrates, novel implantation
    techniques), measure the \Cercis{} novelty component $N(\mathbf{s})$.
    High $N$ should correlate with high inter-laboratory measurement variance,
    validating $N$ as a predictor of certification difficulty.
\end{enumerate}

\subsection{Ablation Study}

To empirically validate the theorems, we specify the following ablation:

\begin{enumerate}[label=\textbf{A\arabic*:}, leftmargin=*]
    \item \textbf{Full audit cohort ($L=10$).} All laboratories active,
    full protocol diversity. Measure $L_{\eff}$ and consensus error rate.
    Compare against Theorem~1 bound.

    \item \textbf{Remove correlated laboratories.} Identify laboratory pairs
    with $\rho_{\ell\ell'} > 0.5$ (shared equipment, shared calibration).
    Drop one from each pair. Expected: $L_{\eff}$ decreases by less than
    $\Delta L$, demonstrating the sub-additivity of correlated information.

    \item \textbf{Reduce protocol diversity ($\cP_\ell = \{\text{Hahn}\}$ for all $\ell$).}
    All laboratories use only Hahn echo. Expected: $L_{\eff}$ drops sharply
    because protocol correlation $\kappa_{pp'} = 1$ for identical protocols,
    inflating $\bar{\rho}$.

    \item \textbf{Vary calibration set size.} Use $n_{\text{cal}} \in \{5, 10, 20, 40\}$.
    Verify that the $\cO(1/\sqrt{n_{\text{cal}}})$ calibration error term
    in Theorem~2 scales as predicted.

    \item \textbf{Test unidentifiability claim.} Provide laboratories with
    only $\Ttwo$ measurements from B2 (depth series) without depth information.
    Ask them to classify each sample as ``surface-limited'' or ``bulk-limited.''
    Theorem~3 predicts classification accuracy at chance level (33\% for
    3-class, 50\% for 2-class), confirming unidentifiability without S1.
\end{enumerate}

% ════════════════════════════════════════════════════════════════════════
\section{Discussion 讨论}
\label{sec:discussion}
% ════════════════════════════════════════════════════════════════════════

\subsection{Summary of Contributions}

This paper establishes the mathematical foundation for auditable quality
certification of nitrogen-vacancy centers in diamond. The three theorems
address complementary aspects of the audit problem:

\begin{enumerate}[nosep]
    \item \textbf{Theorem~1} provides an exponential error detection guarantee:
    the probability that $L_{\eff}$ effectively independent laboratories
    all miss a systematic measurement error decays as $\exp(-2 L_{\eff} \Delta^2)$.
    This formalizes the intuition that multi-laboratory auditing provides
    robustness, while the $L_{\eff}$ correction (accounting for inter-lab
    correlation) prevents overconfidence from nominally large but correlated
    laboratory cohorts.

    \item \textbf{Theorem~2} establishes that the \Cercis{} score is a
    statistically consistent estimator of true \NV{} center quality,
    converging at rate $\cO(\exp(-L_{\eff} \varepsilon^2 / 8))$ in the
    effective laboratory count. This legitimizes the \Cercis{} score as
    more than a heuristic — it is a proper scoring rule for quantum
    device quality.

    \item \textbf{Theorem~3} proves that the causal attribution of $\Ttwo$
    degradation — whether from surface noise, bulk impurities, or measurement
    protocol artifacts — is fundamentally unidentifiable without declared
    structural assumptions. This is the \NV{} center analogue of Derman's
    model risk principle, with direct implications for publication standards.
\end{enumerate}

\subsection{The NV-MIS Publication Standard}

As a direct consequence of Theorem~3, we propose the \textbf{NV Minimum
Information Standard (NV-MIS)} for any publication reporting \NV{} center
coherence properties. Every paper claiming an \NV{} center $\Ttwo$, $\Tone$,
or $\FidReadout$ value must declare:

\begin{enumerate}[label=\textbf{MIS-\arabic*:}, nosep]
    \item \textbf{Depth}: \NV{} implantation energy, simulated or measured
    depth distribution (mean $\pm$ straggle), and depth verification method
    (e.g., NMR depth calibration, ion channeling);
    \item \textbf{Impurities}: Bulk nitrogen concentration (EPR or SIMS),
    ${}^{13}$C isotopic abundance (SIMS or Raman), and other identified
    paramagnetic defect concentrations;
    \item \textbf{Protocol}: Complete pulse sequence specification, microwave
    and laser parameters, repetition rate, data acquisition duration,
    environmental conditions (temperature, magnetic field stability);
    \item \textbf{Surface}: Surface termination chemistry, cleaning procedure,
    and any post-implantation treatment (annealing temperature, duration,
    atmosphere);
    \item \textbf{Substrate}: Diamond growth method (CVD/HPHT), vendor,
    substrate orientation ($\{100\}$, $\{111\}$, $\{110\}$), and any
    pre-treatment.
\end{enumerate}

Journals adopting NV-MIS would eliminate the audit gap identified in
Section~\ref{sec:intro}: a reviewer could, in principle, determine whether
a reported $\Ttwo$ value is consistent with the declared physical parameters
or whether it represents an extraordinary claim requiring extraordinary
multi-laboratory verification.

\subsection{Limitations and Honest Caveats}

We state these limitations explicitly:

\begin{enumerate}[label=\textbf{L\arabic*:}, nosep]
    \item \textbf{Gaussian noise assumption.} The measurement model~\eqref{eq:measurement_model}
    assumes Gaussian noise. Non-Gaussian noise (e.g., telegraph noise from
    fluctuating charges, $1/f$ noise in the electronics) would require
    heavier-tailed concentration inequalities (e.g., Bernstein bounds).
    \item \textbf{Stationary noise spectra.} The analysis assumes time-stationary
    noise spectra $S(\omega)$. Real \NV{} centers experience spectral diffusion
    (slow fluctuations in the noise environment) that violates stationarity
    on timescales of minutes to hours. The theorems' guarantees degrade
    under non-stationarity by a factor proportional to the spectral diffusion
    rate.
    \item \textbf{Finite calibration set.} The $\cO(1/\sqrt{n_{\text{cal}}})$
    term in Theorem~2 becomes non-negligible for small $n_{\text{cal}}$.
    A calibration set of $n_{\text{cal}} \geq 50$ is recommended for
    certification applications.
    \item \textbf{Protocol diversity is bounded.} The four protocols in
    $\cP$ are not fully independent; all share the same optical readout
    stage. A genuinely orthogonal measurement (e.g., electrical readout
    via photocurrent detection) would increase effective information but
    is not widely deployed.
    \item \textbf{No quantum supremacy claims.} This paper makes no claim
    about quantum computational advantage. \NV{} centers are a platform
    for quantum sensing, networking, and information processing; the
    certification framework is orthogonal to the question of whether any
    quantum device outperforms classical computation.
\end{enumerate}

\subsection{Comparison with Existing Approaches}

Current approaches to \NV{} center quality assurance include:

\begin{itemize}[nosep]
    \item \textbf{Single-lab characterization with error bars}: The dominant
    paradigm. Provides no cross-validation against systematic biases.
    \item \textbf{Round-robin inter-laboratory comparisons}: Ad-hoc studies
    (e.g., ``Q-NMR Round Robin'' efforts) that measure the same sample
    across labs but lack a theoretical framework for interpreting disagreement.
    \item \textbf{Standard reference materials}: NIST-traceable diamond
    samples with \NV{} centers of known properties. Valuable but limited
    to specific parameter regimes and not universally available.
\end{itemize}

The \SCX{} framework differs from all three by providing (i)~formal
probabilistic guarantees via Theorems~1--3, (ii)~a scalar audit metric
(\Cercis{} score) that integrates reproducibility and novelty, and
(iii)~a weighting mechanism (\Yajie{} consensus) that adapts to
calibration uncertainty.

\subsection{Future Directions}

\begin{enumerate}[nosep]
    \item \textbf{Extension to other quantum platforms}: The \SCX{} auditing
    framework applies equally to superconducting qubits ($\Ttwo$, gate fidelity),
    trapped ions, and silicon spin qubits. The state space and protocols
    change, but the mathematical structure is invariant.
    \item \textbf{Dynamic auditing}: Continuous monitoring of \NV{} center
    quality over time (aging, radiation damage, surface degradation) within
    the \SCX{} framework.
    \item \textbf{Cryptographic certification}: Integrating the \SCX{} audit
    with blockchain-anchored measurement logs for tamper-proof quality
    certification.
    \item \textbf{Automated protocol selection}: Using the \Spring{} gating
    mechanism to automatically select optimal measurement protocols based
    on real-time noise characterization.
\end{enumerate}

% ════════════════════════════════════════════════════════════════════════
\section{Conclusion 结论}
\label{sec:conclusion}
% ════════════════════════════════════════════════════════════════════════

The nitrogen-vacancy center in diamond is a remarkable quantum system — but its
quality claims are only as credible as the audit framework that certifies them.
This paper has provided that framework.

We have formalized multi-laboratory \NV{} center characterization as an
\SCX{} multi-expert consensus problem, where laboratories are experts,
measurement protocols are modalities, and diamond heterogeneity defines
the state space. Three theorems with full proof ($\rigorFull{}$) establish:

\begin{itemize}[nosep]
    \item Exponentially strong error detection guarantees under multi-laboratory
    auditing (Theorem~1), with the effective laboratory count $L_{\eff}$
    correcting for inter-laboratory correlation from shared resources;
    \item Statistical consistency of the \Cercis{} score as a quality metric,
    converging to true quality at a rate controlled by $L_{\eff}$ and
    calibration set size (Theorem~2);
    \item Fundamental unidentifiability of $\Ttwo$ degradation sources
    (surface noise vs bulk impurities vs protocol artifacts) without
    declared structural assumptions (Theorem~3).
\end{itemize}

The \Cercis{} score $S = Q + \eta N$ and the \Yajie{} multi-laboratory
consensus provide operational tools for quantum device certification.
The NV-MIS publication standard codifies the minimum information required
for auditable \NV{} center quality claims.

\textbf{The central message} is that auditability is not an afterthought
for quantum device characterization — it is a prerequisite for scientific
credibility. A coherence time measured by one laboratory with one protocol
on one diamond sample is not a certified property of the \NV{} center;
it is a claim that requires multi-laboratory verification under the
framework developed here. The mathematics shows that this verification
can be exponentially strong, that the \Cercis{} score provides a principled
quality metric, and that without declaring assumptions, the sources of
discrepancy are fundamentally unidentifiable.

No quantum supremacy is claimed. The supremacy is in the audit.

% ════════════════════════════════════════════════════════════════════════
% REFERENCES
% ════════════════════════════════════════════════════════════════════════
\begin{thebibliography}{99}

\bibitem{doherty2013nitrogen}
M.~W.~Doherty, N.~B.~Manson, P.~Delaney, F.~Jelezko, J.~Wrachtrup, and
L.~C.~L.~Hollenberg, ``The nitrogen-vacancy colour centre in diamond,''
\emph{Physics Reports}, vol.~528, no.~1, pp.~1--45, 2013.

\bibitem{barry2020sensitivity}
J.~F.~Barry, J.~M.~Schloss, E.~Bauch, M.~J.~Turner, C.~A.~Hart,
L.~M.~Pham, and R.~L.~Walsworth, ``Sensitivity optimization for
NV-diamond magnetometry,'' \emph{Reviews of Modern Physics}, vol.~92,
no.~1, p.~015004, 2020.

\bibitem{bradley2019ten}
C.~E.~Bradley, J.~Randall, M.~H.~Abobeih, R.~C.~Berrevoets,
M.~J.~Degen, M.~A.~Bakker, M.~Markham, D.~J.~Twitchen, and
T.~H.~Taminiau, ``A ten-qubit solid-state spin register with quantum
memory up to one minute,'' \emph{Physical Review X}, vol.~9, no.~3,
p.~031045, 2019.

\bibitem{staudacher2013nuclear}
T.~Staudacher, F.~Shi, S.~Pezzagna, J.~Meijer, J.~Du, C.~A.~Meriles,
F.~Reinhard, and J.~Wrachtrup, ``Nuclear magnetic resonance spectroscopy
on a (5-nanometer)$^3$ sample volume,'' \emph{Science}, vol.~339,
no.~6119, pp.~561--563, 2013.

\bibitem{hermans2023entangling}
S.~L.~N.~Hermans, M.~Pompili, H.~K.~C.~Beukers, S.~Baier,
J.~Borregaard, and R.~Hanson, ``Entangling remote qubits using
a single-photon interface,'' \emph{Nature}, vol.~618, pp.~265--270,
2023.

\bibitem{abrahao2025quantum}
R.~Abrahão \emph{et al.}, ``Quantum information processing with
NV centers in diamond,'' \emph{Reviews of Modern Physics}, 2025.

\bibitem{scx2025}
SCX Research Collective, ``The SCX Audit Mandate: Why M-Parameter
Declaration Must Be a Prerequisite for Scientific Publication,''
\emph{SCX Technical Report}, June 2026.

\bibitem{azuma1967}
K.~Azuma, ``Weighted sums of certain dependent random variables,''
\emph{Tohoku Mathematical Journal}, vol.~19, no.~3, pp.~357--367, 1967.

\bibitem{hoeffding1963}
W.~Hoeffding, ``Probability inequalities for sums of bounded random
variables,'' \emph{Journal of the American Statistical Association},
vol.~58, no.~301, pp.~13--30, 1963.

\bibitem{derman1996model}
E.~Derman, ``Model risk,'' \emph{Risk}, vol.~9, no.~5, pp.~34--37, 1996.
Reprinted in \emph{Models.Behaving.Badly}, Free Press, 2011.

\bibitem{balasubramanian2008ultralong}
G.~Balasubramanian \emph{et al.}, ``Ultralong spin coherence time in
isotopically engineered diamond,'' \emph{Nature Materials}, vol.~8,
pp.~383--387, 2009.

\bibitem{bar-gill2013solid}
N.~Bar-Gill, L.~M.~Pham, A.~Jarmola, D.~Budker, and R.~L.~Walsworth,
``Solid-state electronic spin coherence time approaching one second,''
\emph{Nature Communications}, vol.~4, p.~1743, 2013.

\bibitem{myers2014probing}
B.~A.~Myers, A.~Das, M.~C.~Dartiailh, K.~Ohno, D.~D.~Awschalom, and
A.~C.~Bleszynski Jayich, ``Probing surface noise with depth-calibrated
NV centers in diamond,'' \emph{Physical Review Letters}, vol.~113,
no.~2, p.~027602, 2014.

\bibitem{romach2015spectroscopy}
Y.~Romach, C.~Müller, T.~Unden, L.~J.~Rogers, T.~Isoda, K.~M.~Itoh,
M.~Markham, A.~Stacey, J.~Meijer, S.~Pezzagna, B.~Naydenov,
L.~P.~McGuinness, N.~Bar-Gill, and F.~Jelezko, ``Spectroscopy of
surface-induced noise using shallow spins in diamond,''
\emph{Physical Review Letters}, vol.~114, no.~1, p.~017601, 2015.

\bibitem{degen2017quantum}
C.~L.~Degen, F.~Reinhard, and P.~Cappellaro, ``Quantum sensing,''
\emph{Reviews of Modern Physics}, vol.~89, no.~3, p.~035002, 2017.

\bibitem{sangtawesin2019origins}
S.~Sangtawesin, B.~L.~Dwyer, S.~Srinivasan, J.~J.~Allred,
L.~V.~H.~Rodgers, K.~De Greve, A.~Stacey, N.~Dontschuk,
K.~M.~O'Donnell, D.~Hu, D.~A.~Evans, C.~Jayaprakash,
S.~A.~Lyon, and J.~R.~Petta, ``Origins of diamond surface noise
probed by correlating single-spin measurements with surface
spectroscopy,'' \emph{Physical Review X}, vol.~9, no.~3, p.~031052, 2019.

\bibitem{chaudhry2014decoherence}
S.~Chaudhry, ``Decoherence induced by a fluctuating Aharonov-Casher
phase,'' \emph{Physical Review A}, vol.~90, no.~4, p.~042101, 2014.

\bibitem{wang2012protection}
Z.-H.~Wang, G.~de Lange, D.~Ristè, R.~Hanson, and V.~V.~Dobrovitski,
``Comparison of dynamical decoupling protocols for a nitrogen-vacancy
center in diamond,'' \emph{Physical Review B}, vol.~85, no.~15,
p.~155204, 2012.

\bibitem{ma2011efficient}
J.~Ma, X.~Wang, C.~P.~Sun, and F.~Nori, ``Quantum spin squeezing,''
\emph{Physics Reports}, vol.~509, no.~2--3, pp.~89--165, 2011.

\end{thebibliography}

\end{document}
